% ADR-14 — se incluye desde DescripcionArquitecturaSoftware.tex con \input.
% Vive en su propio archivo porque es el ADR más largo del documento y así el
% SAD no crece sin control; el resto de los ADR siguen en línea.

\subsection*{ADR-14 --- Arquitectura híbrida entre SOA y microservicios}

\textbf{Estado:} \textit{supera a ADR-12.} La separación por dominio sigue
vigente; lo que se revisa es el \textbf{criterio} con el que se decide qué
dominio merece ser un despliegue propio, y el resultado pasa de tres servicios
de dominio a \textbf{siete microservicios y siete servicios de dominio}.

\textbf{Contexto:} ADR-12 adoptó una arquitectura basada en servicios con un
criterio de extracción de dos condiciones --- curva de carga e independencia de
disponibilidad --- y produjo tres servicios de dominio. Al recorrer los seis ASR
contra los 32 casos de uso, ese criterio resulta \textbf{insuficiente en dos
sentidos}:

\begin{itemize}
  \item \textbf{No reconoce los servicios transversales.} Pagos lo usan la
        compra (CU-001), la reventa (CU-006), los reembolsos (CU-003), los
        pedidos de alimentos (CU-011) y los parqueaderos (CU-021).
        Autenticación lo usan todas las peticiones del sistema.
        Notificaciones lo usa casi todo. Que uno de ellos falle no degrada un
        dominio: degrada muchos a la vez. ADR-12 no tenía forma de expresar eso.
  \item \textbf{No reconoce el almacén propio.} El Catálogo (CU-005) necesita un
        índice de búsqueda para filtrar por categoría, fecha, ciudad o artista;
        la Cola virtual necesita una estructura de fila en memoria; Acceso QR
        necesita una caché del estado de cada código para validar sin red. Son
        servicios cuyas réplicas \textbf{no competirían} por la base central,
        precisamente porque su estado no vive ahí.
\end{itemize}

Hay además un hecho de carga que ADR-12 sí veía pero resolvía a medias: en la
apertura de venta de un evento muy demandado, la mayor parte del tráfico es de
\textbf{consulta} del catálogo, no de compra. Con ambas cosas en el mismo
servicio, las dos se degradan.

\textbf{Alternativas consideradas:}
\begin{itemize}
  \item \textbf{Todo en microservicios.} Un despliegue por dominio para los 32
        casos de uso. Multiplica bases de datos, despliegues y monitoreo sin
        beneficio para dominios de carga estable --- personal, parqueaderos,
        alimentos ---. Es exactamente el error que ADR-12 corrigió, y con un
        equipo pequeño vuelve a dejar dominios sin implementar.
  \item \textbf{Todo en \textit{service-based}} (lo que fijó ADR-12). Con una
        sola base central y servicios grandes, el pico de la apertura de venta
        obliga a escalar el sistema entero, y las compras compiten por la misma
        base que el resto de los dominios.
  \item \textbf{Híbrida entre SOA y microservicios} --- la elegida.
\end{itemize}

\textbf{Decisión:} la base del sistema es \textit{service-based} ---servicios
grandes por dominio que comparten una base central con un esquema por
servicio--- y de ella se \textbf{extraen como microservicios}, con base de datos
propia y réplicas que escalan solas, los servicios que cumplen al menos
\textbf{dos de cuatro} criterios:

\begin{enumerate}
  \item \textbf{Picos de carga muy por encima del resto:} necesita sumar
        réplicas sin escalar todo el sistema.
  \item \textbf{Lo usan varios dominios:} si se cae, afecta a muchos casos de
        uso a la vez.
  \item \textbf{Su estado cabe en un almacén propio} ---caché, base pequeña,
        índice de búsqueda---: sus réplicas no compiten por la base central.
  \item \textbf{Su falla debe quedar aislada:} un problema en él no puede tumbar
        la operación del evento.
\end{enumerate}

El resultado son \textbf{siete microservicios} ---Cola virtual, Catálogo,
Ventas, Pagos, Autenticación, Acceso QR y Notificaciones--- y \textbf{siete
servicios de dominio} ---Eventos y recintos, Postventa, Alimentos,
Parqueaderos, Personal y logística, Operación del evento, y Finanzas y
reportes---. Todos se exponen por un \textbf{único API Gateway} (ADR-02, que no
cambia) y se comunican de forma asíncrona por un \textbf{broker de eventos}.

\medskip
Tres decisiones acompañan a la principal:

\begin{itemize}
  \item \textbf{Los servicios de dominio también corren redundados, pero no
        escalan.} Dos instancias para disponibilidad, número fijo. Operación del
        evento ---que atiende la evacuación, CU-010--- es crítico para la
        disponibilidad aunque su carga sea baja: necesita redundancia, no
        réplicas elásticas.
  \item \textbf{Finanzas y reportes lee de una réplica de lectura} de la base
        central, no de la principal. Es lo que impide que sus consultas pesadas
        (CU-032) degraden la operación durante el evento, que es el sexto ASR.
  \item \textbf{La compra se resuelve con una Saga orquestada por Ventas}, con
        reserva atómica y vencimiento en la caché, y con llaves de idempotencia
        en Pagos. La compra termina completa o se deshace sin dejar entradas
        bloqueadas; un reintento del comprador o un \textit{webhook} reenviado
        por la pasarela no cobran dos veces (CU-001B, CU-001C, CU-001D).
\end{itemize}

\textbf{Atributo de calidad:} \textbf{Rendimiento} y \textbf{Escalabilidad}
como motivación principal ---los cuatro servicios del camino de compra escalan
sin arrastrar al resto---; \textbf{Disponibilidad}, porque aislar Pagos,
Autenticación y Notificaciones acota el radio de cada falla;
\textbf{Consistencia}, por la Saga y la reserva atómica; y
\textbf{Modificabilidad}, porque los servicios se enteran de los cambios por
eventos y añadir un consumidor nuevo no obliga a tocar a quien los publica.

\textbf{Consecuencias:} es la decisión más cara del documento y conviene no
disimularlo.

\begin{itemize}
  \item \textbf{De 3 unidades desplegables a 14.} La complejidad operativa que
        ADR-12 redujo vuelve a subir, ahora de forma deliberada y acotada: solo
        siete servicios escalan, y los otros siete comparten base y ciclo de
        despliegue.
  \item \textbf{El Núcleo de ADR-12 se deshace.} Aquel agrupó Administración y
        Pedidos en un proceso; aquí Autenticación es microservicio con base
        propia y Alimentos es servicio de dominio sobre la base central. Es el
        único punto donde esta decisión \textit{revierte} trabajo ya hecho, y
        está anotado como tal.
  \item \textbf{Cambia la infraestructura, no solo el código.} El broker pasa a
        uno con tópicos replicados, la orquestación pasa a un clúster con
        autoescalado horizontal, el Catálogo requiere un índice de búsqueda y
        los frontends se sirven por CDN. Nada de eso existe hoy.
  \item \textbf{Riesgo de consistencia eventual.} Catálogo se alimenta por
        eventos de Eventos y recintos, así que un evento recién publicado tarda
        en aparecer en la búsqueda. Es aceptable para CU-005 y no lo es para el
        inventario, que por eso vive en Ventas con reserva atómica.
\end{itemize}

\textbf{Estado de implementación.} A la fecha de este registro la decisión está
\textbf{tomada y no ejecutada}. Lo desplegado sigue siendo lo de ADR-12: tres
servicios de dominio, RabbitMQ, Docker Compose y seis bases separadas. La
distancia entre este diseño y el código está medida en
\texttt{Documentation/Work/MIGRACION-ADR-14.md}, que ordena el camino por fases
y dice qué depende de qué.

\textbf{Cuándo revisar:} si al medir la apertura de venta con carga real alguno
de los cuatro servicios del camino de compra no necesitara escalar de forma
independiente, vuelve a ser candidato a servicio de dominio. El criterio de los
cuatro puntos se aplica igual a cualquier dominio que entre al sistema después.
